\chapter{Approaches to BDD parallelization}
\label{chap:other-impls}

In this chapter, we will go over existing BDD libraries that use parallel execution of the \texttt{apply} operation. We will
explore their design and implementation and assess their strengths and weaknesses. The most prominent BDD libraries
are in Table \ref{tab:bdd-libraries}. All of these libraries are written in low-level programming languages —--
either C or Rust. Only the last two of all the listed libraries implement parallelism: OxiDD and Sylvan. An overview
of both these libraries is presented in the following sections.

\begin{table}[htbp]
  \begin{tabularx}{\textwidth}{llX}
    \toprule
    Library & Language & Parallel? \\
    \midrule
    CUDD              & C    & No  \\
    BuDDy             & C    & No  \\
    biodivine-lib-bdd & Rust & No  \\
    Ruddy             & Rust & No  \\
    OxiDD             & Rust & Yes \\
    Sylvan            & C    & Yes \\
    \bottomrule
  \end{tabularx}
  \caption{Popular BDD libraries}
  \label{tab:bdd-libraries}
\end{table}


\section{OxiDD}
\label{chap:other-impls:oxidd}

OxiDD presents itself as a Decision Diagram framework rather than a library. It focuses on safety, concurrency, and
modularity, and provides powerful abstractions to support its extensibility \cite{oxidd}. The modular architecture
means that the implementation is split into multiple crates\footnote{Crates are the fundamental compilation unit of
Rust code, and can be seen as libraries or projects \cite[keyword.crate]{rust-docs}.}
within the same project. Each crate exposes part of the public API that can be reimplemented to, for example, support
another BDD type - such as MTBDD. In practice, this means you don't add OxiDD as a whole dependency, but rather the
individual OxiDD crates (or their reimplementations).

The most crucial component is the \texttt{core} crate, which consists mainly of trait\footnote{A trait is like an
interface that data types can implement. When a type implements a trait, it can be treated abstractly as that trait
using generics or trait objects \cite[keyword.trait]{rust-docs}.} definitions that formalize the key concepts of DDs
\cite{oxidd}. To use OxiDD as a full-fledged library with a default implementation, use the \texttt{oxidd} crate.

The most significant data structure of the framework is the \texttt{Manager} implemented in the
\texttt{manager-index} and \texttt{manager-pointer} crates. It stores all nodes of a DD and ensures their uniqueness
via a unique table \cite{oxidd}. The difference between these two crates is that \texttt{manager-index} uses 32-bit
unsigned integers to represent edges. The latter offers more flexibility if 32 bits are not enough and nodes are
referenced by pointers. Listing \ref{lst:manager-struct} defines the \texttt{Manager} structure, whose key fields are
examined below.


\begin{listing}[htbp]
    \caption{Definition of the \texttt{Manager} struct in the \texttt{manager-index} crate.}
    \label{lst:manager-struct}
    \begin{minted}{rust}
    #[repr(C)]
    pub struct Manager<'id, DI: DiagramInfo> {
        unique_table: Vec<Mutex<LevelViewSet<'id, DI>>>,
        var_level_map: VarLevelMap,
        var_name_map: VarNameMap,
        data: ManuallyDrop<DI::ManagerData<'id>>,
        store: *const Store<'id, DI>,
        gc_count: AtomicU64,
        reorder_count: u64,
        gc_ongoing: TryLock,
        reorder_gc_prepared: bool,
    }
    \end{minted}
\end{listing}

The \texttt{unique\_table} is composed of several layers of abstraction (aligned with the idea to separate algorithms
from data structures), but essentially boils down to the type \texttt{Vec<Mutex<RawTable>>}. Each item in the vector
represents a level of the BDD. Note that the level does not necessarily refer to a specific variable. OxiDD supports
variable reordering by swapping levels in DDs\footnote{Currently, the \texttt{reorder} crate implements functionality
to establish a given variable order, e.g., harmonize variable orders of different DDs or impose a static variable
order heuristic \cite{oxidd}.}. To maintain synchronization between levels and actual variables, \texttt{Manager}
holds a \texttt{var\_level\_map}. The \texttt{RawTable} is a custom-implemented OxiDD hash table. However, it does not
hold the nodes directly; rather, it holds node IDs --- references to the node store.

The node store contains all nodes of a BDD. Without delving into much detail, the node representation consists of
references to the low and high successors, the level, and a reference counter. To avoid out-of-memory errors, OxiDD
implements garbage collection. The reference counter is used for this exact reason --- only unreferenced nodes are
marked for removal. Garbage collection is an exhaustive operation, so it is delayed as long as possible. It is
executed only if all free space in the node store is depleted, during variable reordering, or manually.

It is important to mention that OxiDD utilizes a \textit{shared representation}. When a single BDD represents
multiple Boolean functions via multiple root handles, it is a shared representation. In a \textit{split
representation}, each function has a separate OBDD, and each graph has a single root \cite{bdds, bfs-bdd, bdds-older}.

Finally, the \texttt{apply} operation. OxiDD uses recursion with a task-based work-stealing architecture. Up to a
certain recursion threshold, tasks are split among the worker threads. After the limit is reached, a single-threaded
algorithm is used \cite{oxidd} --- note that the process remains recursive, but no task splitting for threads is
performed. Because BDDs are rarely balanced, when a thread finishes its shallow subgraph and should become idle, it
tries to steal work from another thread and execute \texttt{apply} on a new subgraph. For this, the \texttt{rayon}
crate is used.

During the execution of the operation, all threads regularly require exclusive access to the same data structures.
These are:

\begin{itemize}
    \item node store
    \item unique table
    \item apply cache
\end{itemize}

To avoid a synchronization bottleneck in the node store, OxiDD maintains it in a shared state across all threads.
Before the operation begins, each thread is given a local portion of the store. Only after the thread fully fills it
is this local store moved to the shared state and receives a new empty slot. The unique table that holds the node IDs
for the node store has a Mutex per BDD level. Lock-free hash tables with atomic operations enable multiple threads to
read and write simultaneously. In this case, the Mutex is used to lock the whole level to grow the table on demand.
This blocking exclusive access is identified as a potential bottleneck for the operation \cite{oxidd}.

The apply cache is the most frequently used data structure during calculation. In the OxiDD implementation, it is a
fixed-size hash table. Table synchronization is performed via a spinlock\footnote{A spinlock is a synchronization
primitive used in concurrent programming. When a thread attempts to acquire a spinlock that is already held by
another thread, it repeatedly executes a loop ("spins") checking if the lock has become available, rather than
yielding control to the operating system scheduler. It is a form of "busy waiting".}. When the lock is held by
another thread, waiting to obtain the lock is too costly, so the result is recomputed instead.


\section{Sylvan}
\label{chap:other-impls:sylvan}
